\section{Discussion}
\label{sec:discussion}

\paragraph{A predictive, qualified answer to the personality question.}
The evidence supports recurring educational behavior in these deployments:
default action distributions learned on one source set improve forecasts on
another, and model-specific student-state responses add information for one
of the three primary actions. Calling this \emph{educational personality} is
useful only with that operational scope. The supported object is a conditional
behavioral profile, with observable forecast error and prompt dependence,
rather than a fixed psychological type assigned to an entire model family.

The three questions constrain each other. RQ1 alone would permit a description
of different output mixtures; RQ2 establishes that part of this information
recurs prospectively. RQ3 shows how instruction changes the conditions under
which those defaults are expressed. Moreover, RQ2 separates a transferable
default from a transferable \emph{response to student state}. The latter is
supported for acknowledgement but remains unresolved for revelation and
elicitation under the tested model class. Combining these outcomes is more
informative than either a general claim that models have different personalities
or a single stability score.

\paragraph{Instruction sensitivity belongs in the description.}
The ASK and EXPLAIN results do not show that an educational personality can
be switched off. They show that the demanded action can become nearly or fully
shared while other permissible actions remain unevenly expressed. For tutor
evaluation, this means that an answer-revelation check alone can miss substantial
differences in what accompanies the answer. It also means that model names or
personality labels are insufficient specifications of a teaching interaction:
the role wording, instructional constraint, student cue, and decoding setting
are part of the empirical description. The neutral-wording results prevent
us from treating the profiles as invariant even under brief role paraphrases.

\paragraph{Increment over existing personality and tutoring research.}
Behavioral personality assessment already studies situated responses and
context-dependent profiles \citep{lee2025trait,liu2026bma}; educational work
already compares tutoring-action mixtures and models tutor-persona variation
\citep{kucheria2025cima,lee2026tutorpersonas}. Our empirical contribution is the
joint finding that prospective default prediction, selective student-state
prediction, and instruction-driven convergence coexist in the same controlled
study. The contrasts distinguish several meanings of recurrence that would
otherwise be collapsed. They do not identify internal persona representations,
establish uniqueness to education, or compare the effectiveness of teaching
styles.

\paragraph{What the archive contributes.}
The existing benchmark supplies breadth and exposes why analysis units matter.
Its tutor, grading, and test-taking roles cannot be pooled into a million
independent observations of a tutor's personality. The source-grouped legacy
analysis also shows how an apparent task-conditioned increment can depend on
an instructional setting. The prospective study responds with matched inputs,
separate problem sources, and explicit forecasts. The archive's scale and the
confirmation experiment's control consequently support different parts of the
argument; neither substitutes for the other.
